\section{Forelesning 2: \emph{Strategisk analyse (del I) og typiske strategiske virkemidler}}
\label{sec:lec02}
\settermlabel{terms:lec02}

\subsection*{Oversikt}
Forelesning 2 introduserer verktøykassen for \emph{ekstern} og \emph{intern}
strategisk analyse, og en katalog av typiske strategiske grep en aktør kan velge.
Målet er å gå fra Porters abstrakte posisjoneringsbegrep (\autoref{sec:lec01})
til konkrete analyseverktøy som strukturerer hva man ser etter når man
diagnostiserer en virksomhets situasjon.

\subsection{Hovedteori}

\subsubsection*{Makroanalyse: \modelref{pestel}{PESTEL} og scenarier}
Før bransjeanalysen kommer makro-omgivelsene. \modelref{pestel}{PESTEL} kartlegger
politiske, økonomiske, sosiale, teknologiske, miljømessige (klima) og
lovgivende faktorer. PESTEL utvides ofte med \emph{forecasting} (megatrender,
svake signaler) og \emph{scenario analysis} -- ulike framtidsbilder testes
mot strategien.

\begin{figure}[H]
\centering
\begin{tikzpicture}[
  every node/.style={font=\sffamily},
  hub/.style={circle, draw=noteBlue, fill=noteBlue, text=white,
              minimum size=1.9cm, font=\sffamily\bfseries\large, inner sep=0pt},
  petal/.style 2 args={circle, draw=#1, fill=#1!12, line width=1pt,
                       minimum size=2.4cm, align=center, inner sep=2pt,
                       font=\sffamily\footnotesize, text=#1!55!black,
                       label={[font=\sffamily\bfseries\Large, text=#1]center:#2}},
]
% Faktorfarger
\definecolor{pestelP}{HTML}{B23A48}
\definecolor{pestelEc}{HTML}{C97B25}
\definecolor{pestelS}{HTML}{D4A017}
\definecolor{pestelT}{HTML}{0E7C7B}
\definecolor{pestelEn}{HTML}{2E7D32}
\definecolor{pestelL}{HTML}{6A1B9A}

% Sentralnav
\node[hub] (c) at (0,0) {PESTEL};

% Seks faktorer plassert i sirkel
\node[petal={pestelP}{P}]  (p)  at ( 90:3.6cm) {\\[10pt]Political};
\node[petal={pestelEc}{E}] (e1) at ( 30:3.6cm) {\\[10pt]Economic};
\node[petal={pestelS}{S}]  (s)  at (330:3.6cm) {\\[10pt]Social};
\node[petal={pestelT}{T}]  (t)  at (270:3.6cm) {\\[10pt]Technological};
\node[petal={pestelEn}{E}] (e2) at (210:3.6cm) {\\[10pt]Environmental};
\node[petal={pestelL}{L}]  (l)  at (150:3.6cm) {\\[10pt]Legal};

% Forbindelseslinjer fra nav til hver faktor
\foreach \n/\col in {p/pestelP, e1/pestelEc, s/pestelS,
                     t/pestelT, e2/pestelEn, l/pestelL}
  \draw[\col, line width=0.9pt, dashed, opacity=0.7] (c) -- (\n);

% Ytre ring som antyder forecasting / scenarier
\draw[noteBlue!40, line width=0.6pt, dotted] (0,0) circle (5.2cm);
\node[noteBlue!70!black, font=\sffamily\itshape\small]
  at (90:5.55cm) {forecasting \& scenarier};
\end{tikzpicture}
\caption{PESTEL-rammeverket: seks makroomgivelser kartlegges rundt et felles
analysenav, og hele kartet leses gjennom forecasting og scenarier
(Johnson et al.).}
\end{figure}

Forecasting-tilnærmingen velges etter hvor stor usikkerheten er: ved lav
usikkerhet kan man bruke et \emph{single-point forecast}, ved moderat
usikkerhet et \emph{interval/range forecast} med sannsynlighetsbånd, og ved
høy usikkerhet \emph{alternative fremtider} (scenarier A/B/C). Tegn på
forandring -- megatrender, inflexion-punkter og svake signaler -- mater inn
i valget av tilnærming.

\begin{figure}[H]
\centering
\includegraphics[width=0.85\linewidth]{lec02-forecasting-under-uncertainty.png}
\caption{Forecasting under usikkerhet: single-point, range og alternative
fremtider -- valg av tilnærming følger usikkerhetsnivået (Johnson et al.\
2017, fig.\ 2.8).}
\end{figure}

\subsubsection*{Bransjeanalyse: \modelref{five-forces}{Porters Five Forces}}
Rammeverket kartlegger konkurransedynamikken i en bransje gjennom fem krefter:
\begin{enumerate}
  \item \term{Industry rivalry} -- intensiteten av konkurranse mellom
    eksisterende aktører.
  \item \term{Threat of new entrants} -- inntrengertrussel, knyttet til
    inngangsbarrierer.
  \item \term{Bargaining power of buyers} -- kjøpermakt.
  \item \term{Bargaining power of suppliers} -- leverandørmakt.
  \item \term{Threat of substitutes} -- substituttrusler.
\end{enumerate}
Den samlede styrken av de fem kreftene bestemmer bransjens strukturelle
lønnsomhetspotensial.

\begin{definition}{Bransjeattraktivitet}
En bransje er strukturelt attraktiv når de fem kreftene er svake. Strategisk
posisjon handler om å plassere seg slik at de fem kreftene treffer en mindre
hardt enn konkurrentene.
\end{definition}

\begin{figure}[H]
\centering
\includegraphics[width=0.85\linewidth]{lec02-porter-five-forces.png}
\caption{Porters Five Forces -- de fem konkurransekreftene som bestemmer
bransjens strukturelle lønnsomhetspotensial.}
\end{figure}

\subsubsection*{\modelref{industry-lifecycle}{Bransjens livssyklus}}
De fem kreftene endrer styrke gjennom bransjens livssyklus: \emph{development
$\rightarrow$ growth $\rightarrow$ shake-out $\rightarrow$ maturity
$\rightarrow$ decline}. F.eks.\ er rivalisering typisk lavest i tidlig
vekstfase og høyest i moden/synkende fase. Modne bransjer (som dansk
dagligvare) har høy rivalisering og lav vekst -- dette har konsekvenser for
hvilke virkemidler som faktisk virker.

\begin{figure}[H]
\centering
\includegraphics[width=0.85\linewidth]{lec02-industry-life-cycle.png}
\caption{Industriens livssyklus med tilhørende endringer i Five Forces
(Johnson et al.).}
\end{figure}

\subsubsection*{Intern analyse: Ressursbasert syn og \modelref{vrio}{VRIO}}
Komplementært til Porters posisjoneringsperspektiv: noen virksomheter har
\modelref{vrio}{VRIO-ressurser} (\emph{Valuable, Rare, Inimitable, Organised}) som gir
\emph{bærekraftig konkurransefortrinn}. VRIO-stigen leder gjennom fire
posisjoner -- fra \emph{vanskelig konkurransesituasjon} til \emph{bærekraftig
konkurransefortrinn}. Relevant for å vurdere f.eks.\ Coops medlemsdata.

\begin{figure}[H]
\centering
\includegraphics[width=0.8\linewidth]{lec02-vrio-decision-matrix.png}
\caption{VRIO-stigen og de fire resulterende konkurranseposisjonene.}
\end{figure}

\subsubsection*{Intern analyse: \modelref{value-chain}{Verdikjeden}}
Porters \modelref{value-chain}{value chain} bryter virksomhetens aktiviteter ned i:
\begin{itemize}
  \item \emph{Primæraktiviteter}: inbound logistics, operations, outbound
    logistics, marketing \& sales, service.
  \item \emph{Støtteaktiviteter}: firm infrastructure, HRM, technology
    development, procurement.
\end{itemize}
Hver aktivitet er en mulig kilde til konkurransefortrinn -- enten gjennom lavere
kostnad eller gjennom differensiering. Verdikjeden er også \emph{baseline} som
Rosenstand \& Baiyere sin \term{Digital Value Chain (DVC)} kontrasteres mot
(jf.\ \autoref{sec:lec03}).

\begin{figure}[H]
\centering
\includegraphics[width=0.85\linewidth]{lec02-porter-value-chain.png}
\caption{Porters klassiske verdikjede.}
\end{figure}

\subsubsection*{\modelref{swot}{SWOT} som syntese}
\modelref{swot}{SWOT} (Styrker, Svakheter, Muligheter, Trusler) syntetiserer
\emph{interne} (S/W) og \emph{eksterne} (O/T) funn fra de foregående
analysene. SWOT er ikke en strategi i seg selv -- den er et hjelpemiddel for
å se sammenhenger.

\begin{figure}[H]
\centering
\includegraphics[width=0.55\linewidth]{lec02-swot.png}
\caption{SWOT-matrisen som syntese av intern og ekstern analyse.}
\end{figure}

\subsubsection*{\modelref{five-building-blocks}{Five Building Blocks of Digital Transformation}}
Ross et al.\ definerer fem grunnleggende kapabiliteter en virksomhet trenger
for å gjennomføre digital transformasjon:
\begin{itemize}
  \item \term{Operational Backbone} -- pålitelige kjerneprosesser og data
    som muliggjør effektiv drift.
  \item \term{Shared Customer Insights} -- felles, oppdatert kundekunnskap
    på tvers av enheter.
  \item \term{Digital Platform} -- gjenbrukbare komponenter for raskt å
    bygge nye digitale tilbud.
  \item \term{Accountability Framework} -- klare beslutningsrettigheter og
    ansvarslinjer for digitale initiativ.
  \item \term{External Developer Platform} -- åpning mot tredjeparter via
    API-er og verktøy.
\end{itemize}
Disse er forutsetninger, ikke strategi -- men strategien må adressere om de
er på plass.

\begin{figure}[H]
\centering
\includegraphics[width=0.85\linewidth]{lec02-five-building-blocks-digital-transformation.png}
\caption{Ross et al.\ sine fem byggeklosser for digital transformasjon.}
\end{figure}

\subsubsection*{Strategy Canvas (Blue Ocean Strategy)}
Kim \& Mauborgne (2005) sin \term{Strategy Canvas} er et diagnostisk og
kreativt verktøy for å visualisere hvordan aktørene i en bransje konkurrerer
-- og hvordan man kan bryte ut av røde havs konkurranse til et \term{blue
ocean}. Canvaset plottes som:
\begin{itemize}
  \item \emph{X-aksen}: bransjens sentrale \term{konkurransefaktorer} (f.eks.\
    pris, sortimentsbredde, butikkdekning, lojalitetsprogram, åpningstid).
  \item \emph{Y-aksen}: nivået kunden opplever på hver faktor (lav $\to$ høy).
  \item \emph{\term{Value curve}} -- den sammenhengende kurven viser aktørens
    strategiske profil og kan sammenlignes med konkurrentenes profiler.
\end{itemize}
For å designe en ny verdikurve brukes \term{Four Actions Framework (ERRC)}:
\begin{itemize}
  \item \term{Eliminate} -- hvilke faktorer som bransjen tar for gitt kan helt fjernes?
  \item \term{Reduce} -- hvilke faktorer bør reduseres godt under bransjestandard?
  \item \term{Raise} -- hvilke faktorer bør heves godt over bransjestandard?
  \item \term{Create} -- hvilke faktorer bør skapes som bransjen aldri har tilbudt?
\end{itemize}
\keypoint{En god strategisk profil har \emph{fokus}, \emph{divergens} fra
bransjens gjennomsnittskurve og et tydelig \emph{slagord} -- canvaset tvinger
fram eksplisitte trade-offs og kobler tilbake til Porters posisjonering
(\autoref{sec:lec01}).}

\subsection{Sentrale fagbegreper}\label{terms:lec02}
\begin{itemize}
  \item \term{PESTEL} -- makro-omgivelsesanalyse.
  \item \term{Five Forces} -- ekstern bransjeanalyse.
  \item \term{Industry life cycle} -- bransjens utviklingsstadier.
  \item \term{Value chain} -- intern aktivitetsanalyse.
  \item \term{VRIO} -- ressursbasert vurdering av varig fortrinn.
  \item \term{SWOT} -- syntese intern $\times$ ekstern.
  \item \term{Entry barriers} -- skalafordeler, brand, byttekostnader, kapital.
  \item \term{Switching costs} -- kunders kostnad ved å bytte leverandør.
  \item \term{Network effects} -- verdi som vokser med antall brukere.
  \item \term{Lock-in} -- mekanisme som binder kunder til en leverandør.
  \item \term{Five Building Blocks} -- digitale grunnkapabiliteter (Ross).
  \item \term{Strategy Canvas} -- visualisering av konkurransefaktorer og
    verdikurve (Kim \& Mauborgne).
  \item \term{Value curve} -- aktørens strategiske profil tegnet på canvaset.
  \item \term{Four Actions Framework (ERRC)} -- Eliminate, Reduce, Raise,
    Create -- verktøy for å designe ny verdikurve.
  \item \term{Blue ocean} -- ukontestert markedsrom skapt ved verdiinnovasjon.
\end{itemize}

\subsection{Typiske strategiske virkemidler}
For en etablert aktør som møter digital disrupsjon (jf.\ \autoref{sec:lec05}):
\begin{enumerate}
  \item \textbf{Forsvare} eksisterende posisjon (styrke eksisterende fortrinn).
  \item \textbf{Investere} i nye digitale kanaler.
  \item \textbf{Kjøpe} digitale kapabiliteter (oppkjøp).
  \item \textbf{Samarbeide} med digitale aktører (allianser, plattformer).
  \item \textbf{Selv-disrupsjon} -- bygge en disruptiv enhet før andre gjør det.
\end{enumerate}
Valget avhenger av strategisk fit med forretningsmodellen (\autoref{sec:lec03}).

\subsection{Eksempler og caser}
\begin{eksempelboks}{Dansk dagligvare -- Five Forces}
Høy rivalisering (Salling, Coop, REMA, Lidl), moderat trussel fra netthandel
(Nemlig), sterk leverandørmakt fra store merker, lav kjøpermakt per kunde men
sterk samlet, substitutter inkluderer ferdigmat-bud og restaurant. Moden
bransje med strukturelt lav margin.
\end{eksempelboks}

\begin{eksempelboks}{Coops verdikjede}
Innkjøp og logistikk er sentralt; digital teknologi i butikk (Scan \& Betal,
lojalitetsapp) hører hjemme under \emph{technology development} og
\emph{marketing \& sales}. Medlemsdataene er VRIO-kandidat -- verdifulle,
relativt sjeldne, vanskelig å imitere uten samme medlemsbase, men
organisasjonen må være riktig satt opp for å utnytte dem.
\end{eksempelboks}

\subsection{Eksamensrelevans}
\begin{itemize}
  \item Kunne gjennomføre en Five Forces-analyse på en konkret bransje.
  \item Kunne plassere digitale initiativ i verdikjeden og forklare hvor
    fortrinn oppstår.
  \item Kjenne katalogen av strategiske virkemidler og kunne velge mellom dem.
  \item Forstå at ekstern + intern analyse er to nødvendige sider av samme sak,
    og at SWOT \emph{syntetiserer}.
  \item Bruke VRIO til å vurdere om en ressurs gir varig fortrinn -- ikke
    bare et øyeblikkelig forsprang.
\end{itemize}
